\section{Related Works}
\label{appendix:related_work}
\paragraph{Discrete diffusion and flows.} Early diffusion models were formulated as continuous-time Markov chains over continuous spaces with Gaussian transition kernels \citep{sohl2015deep,ho2020denoising}, and were later connected to continuous-time formulations via stochastic differential equations, offering a unifying perspective on score-based generative modeling \citep{song2020score}. In parallel, \emph{discrete} diffusion has been developed from the viewpoint of Markov chains over discrete space \citep{hoogeboom2021argmax}. Notably, \citet{austin2021structured} introduced D3PM with several families of discrete transition kernels, and \citet{lou2023discrete} proposed SEDD, which adopts score-based training objectives. A complementary line of work studies \emph{discrete flows} \citep{campbell2024generative,gat2024discrete}, aiming to understand continuous-time Markov chains (CTMCs) that interpolate between data and base distributions; this perspective aligns with ours. Subsequent extensions consider token-wise paths and path-wise structure within such flows \citep{shaul2024flow}.

\paragraph{Masked Diffusion Models.}
Among discrete-transition designs, absorbing-state (a.k.a.\ masking) kernels have become a popular and strong-performing choice. Recent work shows that this yields a simple and principled training recipe, referred to as Masked Diffusion Models (MDMs) \citep{sahoo2024simple,shi2024simplified}. A growing body of results demonstrates the scalability of this approach across problem settings and modalities, including large-scale natural language modeling \citep{nie2024scaling,nie2025large,dream2025,song2025seed,gemini2025diffusion}, code generation \citep{labs2025mercury,gong2025diffucoder}, and multimodal learning \citep{swerdlow2025unified}.

\paragraph{Any-order inference in MDMs.}
With the advent of MDMs, subsequent work has established that they admit theoretically grounded \emph{any-order inference}, wherein tokens can be unmasked in arbitrary orders rather than following a fixed CTMC schedule \citep{kim2025train,peng2025path}. Practical token-ordering rules span a spectrum of heuristics based on model confidence and uncertainty—e.g., maximum-probability logits \citep{chang2022maskgit,zheng2024masked}, probability margin \citep{kim2025train}, semi-autoregressive schedules \citep{nie2024scaling}, and entropy-based criteria \citep{ben2025accelerated}—as well as strategies that leverage reference models to guide the unmasking trajectory \citep{peng2025path}. Beyond heuristics, another thread trains auxiliary modules to anchor or adapt the generation order \citep{rout2025anchored}, while recent work directly \emph{learns} token orders end-to-end \citep{ma2025reinforced,wang2025learning}.


\paragraph{Stochastic interpolant.} Stochastic interpolant ~\citep{albergo2022building, albergo2023stochastic} is a general framework for building measure-transport based generative models on continuous state space. While building off different philosophical grounds, it can be seen as equivalent to flow matching~\citep{lipman2022flow, liu2022flowstraightfastlearning}. Extensions of the interpolant have been proposed for conditional generation through data-dependent coupling ~\citep{albergo2024stochastic}, which we adopt for infilling task design in Section~\ref{sec:experiment}.


\paragraph{Descriptive overview on concurrent work.} The most notable concurrent work is EditFlow~\citep{havasi2025edit}, where the primary mathematical machinery that enabled their construction is referred to as ``Flow Matching with Auxiliary Process". We note that this can be interpreted as mathematically equivalent to the notion of joint interpolant in this work, e.g., our Proposition~\ref{prop:joint-target-rate} is equivalent to Theorem 3.1 in \cite{havasi2025edit}.

The main differences are (1) the choice of interpolant and (2) the guarantee of any-order inference. Whereas EditFlow is built around an explicit probability path, we instead define a pair of coupled random variables that implicitly induce this path, leading to a different choice of intermediate. As discussed in Section~\ref{sec:FlexMDM_informal}, our choice of interpolant yields a distinct training objective for an unmasking posterior and an insertion-expectation term. Consequently, it enables the any-order inference guarantee established in Section~\ref{sec:FlexMDM_inference}. 